% Neuron document style: Charter + Noto Serif CJK KR, light outlined boxes, grayscale.
% Usage: \documentclass[11pt]{report} % Neuron document style: Charter + Noto Serif CJK KR, light outlined boxes, grayscale.
% Usage: \documentclass[11pt]{report} % Neuron document style: Charter + Noto Serif CJK KR, light outlined boxes, grayscale.
% Usage: \documentclass[11pt]{report} % Neuron document style: Charter + Noto Serif CJK KR, light outlined boxes, grayscale.
% Usage: \documentclass[11pt]{report} \input{neuron-style}  (compile with xelatex)
\usepackage[a4paper,margin=24mm,headheight=14pt]{geometry}
\usepackage{fontspec}
\setmainfont{Charter}[Path=fonts/,Extension=.otf,UprightFont=*-Regular,ItalicFont=*-Italic,
  BoldFont=*-Bold,BoldItalicFont=*-BoldItalic,Ligatures=TeX]
\setsansfont{Latin Modern Sans}
\setmonofont{Latin Modern Mono}
\newfontfamily\krfont{Noto Serif CJK KR}[Scale=0.9]
\XeTeXinterchartokenstate=1
\newXeTeXintercharclass\krclass
\count255="AC00 \loop\ifnum\count255<"D7A4 \XeTeXcharclass\count255=\krclass \advance\count255 by 1 \repeat
\count255="3000 \loop\ifnum\count255<"3040 \XeTeXcharclass\count255=\krclass \advance\count255 by 1 \repeat
\count255="3130 \loop\ifnum\count255<"3190 \XeTeXcharclass\count255=\krclass \advance\count255 by 1 \repeat
\count255="1100 \loop\ifnum\count255<"1200 \XeTeXcharclass\count255=\krclass \advance\count255 by 1 \repeat
\XeTeXinterchartoks 0 \krclass = {\begingroup\krfont}
\XeTeXinterchartoks 4095 \krclass = {\begingroup\krfont}
\XeTeXinterchartoks \krclass 0 = {\endgroup}
\XeTeXinterchartoks \krclass 4095 = {\endgroup}
\XeTeXlinebreaklocale "ko"
\XeTeXlinebreakskip=0pt plus 1pt
\usepackage{amsmath,amssymb,amsthm,mathtools,bm}
\usepackage[italic,defaultmathsizes,noasterisk]{mathastext}
\emergencystretch=2.5em

\usepackage{graphicx,booktabs,tabularx,longtable,array,multirow,enumitem,xcolor}
\usepackage{caption}
\captionsetup{font=small,labelfont=bf}
\usepackage{tikz}
\usetikzlibrary{shapes.geometric,arrows.meta,positioning,fit,calc,backgrounds,decorations.pathreplacing}
\usepackage[most]{tcolorbox}
\usepackage{listings}
\usepackage{fancyhdr}
\usepackage[hidelinks]{hyperref}

\setlist{itemsep=1pt,topsep=3pt,parsep=0pt}
\setlength{\parskip}{3pt}
\renewcommand{\arraystretch}{1.12}

\lstdefinestyle{sql}{language=SQL,basicstyle=\ttfamily\small,keywordstyle=\bfseries,
  commentstyle=\itshape\color{black!60},columns=fullflexible,keepspaces=true,
  frame=l,framerule=0.6pt,rulecolor=\color{black!40},xleftmargin=8pt,framexleftmargin=4pt,
  aboveskip=5pt,belowskip=5pt,showstringspaces=false,breaklines=true,
  morekeywords={OVER,PARTITION,RANK,DENSE_RANK,ROW_NUMBER,NTILE,LAG,LEAD,ROWS,RANGE,PRECEDING,FOLLOWING,UNBOUNDED,CURRENT,ROW,NVL,NVL2,REFERENCES,CASCADE,MINUS,INTERSECT,EXISTS,MATERIALIZED,VIEW,REPLACE,RATIO_TO_REPORT,FIRST_VALUE,LAST_VALUE,INTERVAL,NULLS,LAST,FIRST,CHECK,CONSTRAINT,UNIQUE,DEFERRABLE,COMMIT,BEGIN,ROLLBACK,TRANSACTION,INDEX,DEFAULT,ANY,ALL,SOME,ROUND}}
\lstset{style=sql}

\tcbset{
 nb/.style={enhanced,breakable,colback=white,coltitle=black,fonttitle=\bfseries,
   left=3mm,right=3mm,top=1.6mm,bottom=1.6mm,before skip=9pt,after skip=9pt,parbox=false},
 framed/.style={nb,boxrule=0.5pt,colframe=black!70,arc=1.2mm,colbacktitle=white,
   titlerule=0.4pt,titlerule style={black!25},toptitle=1mm,bottomtitle=0.6mm},
 light/.style={framed,colframe=black!35},
 ruled/.style={nb,boxrule=0pt,leftrule=1.1pt,colframe=black!45,sharp corners,left=3.5mm,
   top=0.5mm,bottom=0.5mm,detach title,before upper={\tcbtitle\par\smallskip}},
}
\DeclareTColorBox{keybox}{o}{framed,leftrule=2.2pt,title={\IfValueTF{#1}{#1}{Key idea}}}
\DeclareTColorBox{formulabox}{o}{framed,IfValueTF={#1}{title={#1}}{}}
\DeclareTColorBox{workout}{o}{light,title={Example\IfValueT{#1}{\ \textmd{---}\ #1}}}
\DeclareTColorBox{qa}{m}{light,title={#1}}
\DeclareTColorBox{warnbox}{o}{ruled,leftrule=2.2pt,colframe=black,title={Caution\IfValueT{#1}{\ \textmd{---}\ #1}}}
\DeclareTColorBox{whybox}{o}{ruled,title={\IfValueTF{#1}{#1}{Why it works}}}
\DeclareTColorBox{prof}{o}{ruled,title={Lecture note\IfValueT{#1}{\ \textmd{---}\ #1}}}
\DeclareTColorBox{career}{o}{ruled,title={Career note\IfValueT{#1}{\ \textmd{---}\ #1}}}
\DeclareTColorBox{insight}{o}{ruled,title={In practice\IfValueT{#1}{\ \textmd{---}\ #1}}}

\newcommand{\Step}[1]{\par\smallskip\noindent\textbf{Step #1.}\ }
\newcommand{\Concept}[1]{\hfill{\small[개념: #1]}}
\newcommand{\Use}[1]{\noindent\textit{사용 개념:} #1\par}
\newcommand{\Ans}[1]{\par\smallskip\noindent\fbox{\parbox{\dimexpr\linewidth-2\fboxsep-2\fboxrule\relax}{\textbf{답}\ \ #1}}\par}
\newcommand{\Trap}[1]{\par\smallskip\noindent\textit{흔한 실수:} #1}
\newcommand{\Verify}[1]{\par\smallskip\noindent\textit{검산:} #1}
\newcommand{\Intuit}[1]{\par\smallskip\textbf{직관: #1}}
\newcommand{\pk}[1]{\underline{#1}}
\newcommand{\NULL}{\textsc{null}}
\newcommand{\sqlk}[1]{\texttt{#1}}
\newcommand{\Nul}{{\color{black!50}\textit{null}}}

% relational algebra symbols
\newcommand{\sel}{\sigma}
\newcommand{\proj}{\pi}
\newcommand{\join}{\bowtie}

\pagestyle{fancy}
\fancyhf{}
\fancyhead[R]{\small\itshape\leftmark}
\fancyfoot[C]{\small\thepage}
\renewcommand{\headrulewidth}{0pt}
\renewcommand{\chaptermark}[1]{\markboth{#1}{}}

\usepackage{titlesec}
\titleformat{\chapter}[display]{\bfseries\Large}{\chaptertitlename\ \thechapter}{6pt}{\LARGE}
\titlespacing*{\chapter}{0pt}{-10pt}{14pt}
\titleformat{\section}{\bfseries\large}{\thesection}{0.8em}{}
\titleformat{\subsection}{\bfseries\normalsize}{\thesubsection}{0.8em}{}

% ---------- cheat-sheet / problem-sheet helpers ----------
\newcounter{qn}
\newcommand{\Q}[2]{\par\smallskip\refstepcounter{qn}\noindent\textbf{Q\theqn.}\ {\footnotesize[#1]}\ #2\par}
\newcommand{\A}[1]{\vspace{1pt}\noindent\begin{tcolorbox}[ruled,leftrule=1.6pt,colframe=black!45,before skip=1pt,after skip=4pt,top=0.5mm,bottom=0.5mm]\small\textbf{답} #1\end{tcolorbox}}
\newcommand{\Lab}[3]{\item \textbf{#1}\\ \textit{해 본 것:} #2\\ \textit{\ensuremath{\Rightarrow} 개념:} #3}
\newcommand{\eg}{\textit{예)}\ }
\newcommand{\bi}{\textit{비유)}\ }
  (compile with xelatex)
\usepackage[a4paper,margin=24mm,headheight=14pt]{geometry}
\usepackage{fontspec}
\setmainfont{Charter}[Path=fonts/,Extension=.otf,UprightFont=*-Regular,ItalicFont=*-Italic,
  BoldFont=*-Bold,BoldItalicFont=*-BoldItalic,Ligatures=TeX]
\setsansfont{Latin Modern Sans}
\setmonofont{Latin Modern Mono}
\newfontfamily\krfont{Noto Serif CJK KR}[Scale=0.9]
\XeTeXinterchartokenstate=1
\newXeTeXintercharclass\krclass
\count255="AC00 \loop\ifnum\count255<"D7A4 \XeTeXcharclass\count255=\krclass \advance\count255 by 1 \repeat
\count255="3000 \loop\ifnum\count255<"3040 \XeTeXcharclass\count255=\krclass \advance\count255 by 1 \repeat
\count255="3130 \loop\ifnum\count255<"3190 \XeTeXcharclass\count255=\krclass \advance\count255 by 1 \repeat
\count255="1100 \loop\ifnum\count255<"1200 \XeTeXcharclass\count255=\krclass \advance\count255 by 1 \repeat
\XeTeXinterchartoks 0 \krclass = {\begingroup\krfont}
\XeTeXinterchartoks 4095 \krclass = {\begingroup\krfont}
\XeTeXinterchartoks \krclass 0 = {\endgroup}
\XeTeXinterchartoks \krclass 4095 = {\endgroup}
\XeTeXlinebreaklocale "ko"
\XeTeXlinebreakskip=0pt plus 1pt
\usepackage{amsmath,amssymb,amsthm,mathtools,bm}
\usepackage[italic,defaultmathsizes,noasterisk]{mathastext}
\emergencystretch=2.5em

\usepackage{graphicx,booktabs,tabularx,longtable,array,multirow,enumitem,xcolor}
\usepackage{caption}
\captionsetup{font=small,labelfont=bf}
\usepackage{tikz}
\usetikzlibrary{shapes.geometric,arrows.meta,positioning,fit,calc,backgrounds,decorations.pathreplacing}
\usepackage[most]{tcolorbox}
\usepackage{listings}
\usepackage{fancyhdr}
\usepackage[hidelinks]{hyperref}

\setlist{itemsep=1pt,topsep=3pt,parsep=0pt}
\setlength{\parskip}{3pt}
\renewcommand{\arraystretch}{1.12}

\lstdefinestyle{sql}{language=SQL,basicstyle=\ttfamily\small,keywordstyle=\bfseries,
  commentstyle=\itshape\color{black!60},columns=fullflexible,keepspaces=true,
  frame=l,framerule=0.6pt,rulecolor=\color{black!40},xleftmargin=8pt,framexleftmargin=4pt,
  aboveskip=5pt,belowskip=5pt,showstringspaces=false,breaklines=true,
  morekeywords={OVER,PARTITION,RANK,DENSE_RANK,ROW_NUMBER,NTILE,LAG,LEAD,ROWS,RANGE,PRECEDING,FOLLOWING,UNBOUNDED,CURRENT,ROW,NVL,NVL2,REFERENCES,CASCADE,MINUS,INTERSECT,EXISTS,MATERIALIZED,VIEW,REPLACE,RATIO_TO_REPORT,FIRST_VALUE,LAST_VALUE,INTERVAL,NULLS,LAST,FIRST,CHECK,CONSTRAINT,UNIQUE,DEFERRABLE,COMMIT,BEGIN,ROLLBACK,TRANSACTION,INDEX,DEFAULT,ANY,ALL,SOME,ROUND}}
\lstset{style=sql}

\tcbset{
 nb/.style={enhanced,breakable,colback=white,coltitle=black,fonttitle=\bfseries,
   left=3mm,right=3mm,top=1.6mm,bottom=1.6mm,before skip=9pt,after skip=9pt,parbox=false},
 framed/.style={nb,boxrule=0.5pt,colframe=black!70,arc=1.2mm,colbacktitle=white,
   titlerule=0.4pt,titlerule style={black!25},toptitle=1mm,bottomtitle=0.6mm},
 light/.style={framed,colframe=black!35},
 ruled/.style={nb,boxrule=0pt,leftrule=1.1pt,colframe=black!45,sharp corners,left=3.5mm,
   top=0.5mm,bottom=0.5mm,detach title,before upper={\tcbtitle\par\smallskip}},
}
\DeclareTColorBox{keybox}{o}{framed,leftrule=2.2pt,title={\IfValueTF{#1}{#1}{Key idea}}}
\DeclareTColorBox{formulabox}{o}{framed,IfValueTF={#1}{title={#1}}{}}
\DeclareTColorBox{workout}{o}{light,title={Example\IfValueT{#1}{\ \textmd{---}\ #1}}}
\DeclareTColorBox{qa}{m}{light,title={#1}}
\DeclareTColorBox{warnbox}{o}{ruled,leftrule=2.2pt,colframe=black,title={Caution\IfValueT{#1}{\ \textmd{---}\ #1}}}
\DeclareTColorBox{whybox}{o}{ruled,title={\IfValueTF{#1}{#1}{Why it works}}}
\DeclareTColorBox{prof}{o}{ruled,title={Lecture note\IfValueT{#1}{\ \textmd{---}\ #1}}}
\DeclareTColorBox{career}{o}{ruled,title={Career note\IfValueT{#1}{\ \textmd{---}\ #1}}}
\DeclareTColorBox{insight}{o}{ruled,title={In practice\IfValueT{#1}{\ \textmd{---}\ #1}}}

\newcommand{\Step}[1]{\par\smallskip\noindent\textbf{Step #1.}\ }
\newcommand{\Concept}[1]{\hfill{\small[개념: #1]}}
\newcommand{\Use}[1]{\noindent\textit{사용 개념:} #1\par}
\newcommand{\Ans}[1]{\par\smallskip\noindent\fbox{\parbox{\dimexpr\linewidth-2\fboxsep-2\fboxrule\relax}{\textbf{답}\ \ #1}}\par}
\newcommand{\Trap}[1]{\par\smallskip\noindent\textit{흔한 실수:} #1}
\newcommand{\Verify}[1]{\par\smallskip\noindent\textit{검산:} #1}
\newcommand{\Intuit}[1]{\par\smallskip\textbf{직관: #1}}
\newcommand{\pk}[1]{\underline{#1}}
\newcommand{\NULL}{\textsc{null}}
\newcommand{\sqlk}[1]{\texttt{#1}}
\newcommand{\Nul}{{\color{black!50}\textit{null}}}

% relational algebra symbols
\newcommand{\sel}{\sigma}
\newcommand{\proj}{\pi}
\newcommand{\join}{\bowtie}

\pagestyle{fancy}
\fancyhf{}
\fancyhead[R]{\small\itshape\leftmark}
\fancyfoot[C]{\small\thepage}
\renewcommand{\headrulewidth}{0pt}
\renewcommand{\chaptermark}[1]{\markboth{#1}{}}

\usepackage{titlesec}
\titleformat{\chapter}[display]{\bfseries\Large}{\chaptertitlename\ \thechapter}{6pt}{\LARGE}
\titlespacing*{\chapter}{0pt}{-10pt}{14pt}
\titleformat{\section}{\bfseries\large}{\thesection}{0.8em}{}
\titleformat{\subsection}{\bfseries\normalsize}{\thesubsection}{0.8em}{}

% ---------- cheat-sheet / problem-sheet helpers ----------
\newcounter{qn}
\newcommand{\Q}[2]{\par\smallskip\refstepcounter{qn}\noindent\textbf{Q\theqn.}\ {\footnotesize[#1]}\ #2\par}
\newcommand{\A}[1]{\vspace{1pt}\noindent\begin{tcolorbox}[ruled,leftrule=1.6pt,colframe=black!45,before skip=1pt,after skip=4pt,top=0.5mm,bottom=0.5mm]\small\textbf{답} #1\end{tcolorbox}}
\newcommand{\Lab}[3]{\item \textbf{#1}\\ \textit{해 본 것:} #2\\ \textit{\ensuremath{\Rightarrow} 개념:} #3}
\newcommand{\eg}{\textit{예)}\ }
\newcommand{\bi}{\textit{비유)}\ }
  (compile with xelatex)
\usepackage[a4paper,margin=24mm,headheight=14pt]{geometry}
\usepackage{fontspec}
\setmainfont{Charter}[Path=fonts/,Extension=.otf,UprightFont=*-Regular,ItalicFont=*-Italic,
  BoldFont=*-Bold,BoldItalicFont=*-BoldItalic,Ligatures=TeX]
\setsansfont{Latin Modern Sans}
\setmonofont{Latin Modern Mono}
\newfontfamily\krfont{Noto Serif CJK KR}[Scale=0.9]
\XeTeXinterchartokenstate=1
\newXeTeXintercharclass\krclass
\count255="AC00 \loop\ifnum\count255<"D7A4 \XeTeXcharclass\count255=\krclass \advance\count255 by 1 \repeat
\count255="3000 \loop\ifnum\count255<"3040 \XeTeXcharclass\count255=\krclass \advance\count255 by 1 \repeat
\count255="3130 \loop\ifnum\count255<"3190 \XeTeXcharclass\count255=\krclass \advance\count255 by 1 \repeat
\count255="1100 \loop\ifnum\count255<"1200 \XeTeXcharclass\count255=\krclass \advance\count255 by 1 \repeat
\XeTeXinterchartoks 0 \krclass = {\begingroup\krfont}
\XeTeXinterchartoks 4095 \krclass = {\begingroup\krfont}
\XeTeXinterchartoks \krclass 0 = {\endgroup}
\XeTeXinterchartoks \krclass 4095 = {\endgroup}
\XeTeXlinebreaklocale "ko"
\XeTeXlinebreakskip=0pt plus 1pt
\usepackage{amsmath,amssymb,amsthm,mathtools,bm}
\usepackage[italic,defaultmathsizes,noasterisk]{mathastext}
\emergencystretch=2.5em

\usepackage{graphicx,booktabs,tabularx,longtable,array,multirow,enumitem,xcolor}
\usepackage{caption}
\captionsetup{font=small,labelfont=bf}
\usepackage{tikz}
\usetikzlibrary{shapes.geometric,arrows.meta,positioning,fit,calc,backgrounds,decorations.pathreplacing}
\usepackage[most]{tcolorbox}
\usepackage{listings}
\usepackage{fancyhdr}
\usepackage[hidelinks]{hyperref}

\setlist{itemsep=1pt,topsep=3pt,parsep=0pt}
\setlength{\parskip}{3pt}
\renewcommand{\arraystretch}{1.12}

\lstdefinestyle{sql}{language=SQL,basicstyle=\ttfamily\small,keywordstyle=\bfseries,
  commentstyle=\itshape\color{black!60},columns=fullflexible,keepspaces=true,
  frame=l,framerule=0.6pt,rulecolor=\color{black!40},xleftmargin=8pt,framexleftmargin=4pt,
  aboveskip=5pt,belowskip=5pt,showstringspaces=false,breaklines=true,
  morekeywords={OVER,PARTITION,RANK,DENSE_RANK,ROW_NUMBER,NTILE,LAG,LEAD,ROWS,RANGE,PRECEDING,FOLLOWING,UNBOUNDED,CURRENT,ROW,NVL,NVL2,REFERENCES,CASCADE,MINUS,INTERSECT,EXISTS,MATERIALIZED,VIEW,REPLACE,RATIO_TO_REPORT,FIRST_VALUE,LAST_VALUE,INTERVAL,NULLS,LAST,FIRST,CHECK,CONSTRAINT,UNIQUE,DEFERRABLE,COMMIT,BEGIN,ROLLBACK,TRANSACTION,INDEX,DEFAULT,ANY,ALL,SOME,ROUND}}
\lstset{style=sql}

\tcbset{
 nb/.style={enhanced,breakable,colback=white,coltitle=black,fonttitle=\bfseries,
   left=3mm,right=3mm,top=1.6mm,bottom=1.6mm,before skip=9pt,after skip=9pt,parbox=false},
 framed/.style={nb,boxrule=0.5pt,colframe=black!70,arc=1.2mm,colbacktitle=white,
   titlerule=0.4pt,titlerule style={black!25},toptitle=1mm,bottomtitle=0.6mm},
 light/.style={framed,colframe=black!35},
 ruled/.style={nb,boxrule=0pt,leftrule=1.1pt,colframe=black!45,sharp corners,left=3.5mm,
   top=0.5mm,bottom=0.5mm,detach title,before upper={\tcbtitle\par\smallskip}},
}
\DeclareTColorBox{keybox}{o}{framed,leftrule=2.2pt,title={\IfValueTF{#1}{#1}{Key idea}}}
\DeclareTColorBox{formulabox}{o}{framed,IfValueTF={#1}{title={#1}}{}}
\DeclareTColorBox{workout}{o}{light,title={Example\IfValueT{#1}{\ \textmd{---}\ #1}}}
\DeclareTColorBox{qa}{m}{light,title={#1}}
\DeclareTColorBox{warnbox}{o}{ruled,leftrule=2.2pt,colframe=black,title={Caution\IfValueT{#1}{\ \textmd{---}\ #1}}}
\DeclareTColorBox{whybox}{o}{ruled,title={\IfValueTF{#1}{#1}{Why it works}}}
\DeclareTColorBox{prof}{o}{ruled,title={Lecture note\IfValueT{#1}{\ \textmd{---}\ #1}}}
\DeclareTColorBox{career}{o}{ruled,title={Career note\IfValueT{#1}{\ \textmd{---}\ #1}}}
\DeclareTColorBox{insight}{o}{ruled,title={In practice\IfValueT{#1}{\ \textmd{---}\ #1}}}

\newcommand{\Step}[1]{\par\smallskip\noindent\textbf{Step #1.}\ }
\newcommand{\Concept}[1]{\hfill{\small[개념: #1]}}
\newcommand{\Use}[1]{\noindent\textit{사용 개념:} #1\par}
\newcommand{\Ans}[1]{\par\smallskip\noindent\fbox{\parbox{\dimexpr\linewidth-2\fboxsep-2\fboxrule\relax}{\textbf{답}\ \ #1}}\par}
\newcommand{\Trap}[1]{\par\smallskip\noindent\textit{흔한 실수:} #1}
\newcommand{\Verify}[1]{\par\smallskip\noindent\textit{검산:} #1}
\newcommand{\Intuit}[1]{\par\smallskip\textbf{직관: #1}}
\newcommand{\pk}[1]{\underline{#1}}
\newcommand{\NULL}{\textsc{null}}
\newcommand{\sqlk}[1]{\texttt{#1}}
\newcommand{\Nul}{{\color{black!50}\textit{null}}}

% relational algebra symbols
\newcommand{\sel}{\sigma}
\newcommand{\proj}{\pi}
\newcommand{\join}{\bowtie}

\pagestyle{fancy}
\fancyhf{}
\fancyhead[R]{\small\itshape\leftmark}
\fancyfoot[C]{\small\thepage}
\renewcommand{\headrulewidth}{0pt}
\renewcommand{\chaptermark}[1]{\markboth{#1}{}}

\usepackage{titlesec}
\titleformat{\chapter}[display]{\bfseries\Large}{\chaptertitlename\ \thechapter}{6pt}{\LARGE}
\titlespacing*{\chapter}{0pt}{-10pt}{14pt}
\titleformat{\section}{\bfseries\large}{\thesection}{0.8em}{}
\titleformat{\subsection}{\bfseries\normalsize}{\thesubsection}{0.8em}{}

% ---------- cheat-sheet / problem-sheet helpers ----------
\newcounter{qn}
\newcommand{\Q}[2]{\par\smallskip\refstepcounter{qn}\noindent\textbf{Q\theqn.}\ {\footnotesize[#1]}\ #2\par}
\newcommand{\A}[1]{\vspace{1pt}\noindent\begin{tcolorbox}[ruled,leftrule=1.6pt,colframe=black!45,before skip=1pt,after skip=4pt,top=0.5mm,bottom=0.5mm]\small\textbf{답} #1\end{tcolorbox}}
\newcommand{\Lab}[3]{\item \textbf{#1}\\ \textit{해 본 것:} #2\\ \textit{\ensuremath{\Rightarrow} 개념:} #3}
\newcommand{\eg}{\textit{예)}\ }
\newcommand{\bi}{\textit{비유)}\ }
  (compile with xelatex)
\usepackage[a4paper,margin=24mm,headheight=14pt]{geometry}
\usepackage{fontspec}
\setmainfont{Charter}[Path=fonts/,Extension=.otf,UprightFont=*-Regular,ItalicFont=*-Italic,
  BoldFont=*-Bold,BoldItalicFont=*-BoldItalic,Ligatures=TeX]
\setsansfont{Latin Modern Sans}
\setmonofont{Latin Modern Mono}
\newfontfamily\krfont{Noto Serif CJK KR}[Scale=0.9]
\XeTeXinterchartokenstate=1
\newXeTeXintercharclass\krclass
\count255="AC00 \loop\ifnum\count255<"D7A4 \XeTeXcharclass\count255=\krclass \advance\count255 by 1 \repeat
\count255="3000 \loop\ifnum\count255<"3040 \XeTeXcharclass\count255=\krclass \advance\count255 by 1 \repeat
\count255="3130 \loop\ifnum\count255<"3190 \XeTeXcharclass\count255=\krclass \advance\count255 by 1 \repeat
\count255="1100 \loop\ifnum\count255<"1200 \XeTeXcharclass\count255=\krclass \advance\count255 by 1 \repeat
\XeTeXinterchartoks 0 \krclass = {\begingroup\krfont}
\XeTeXinterchartoks 4095 \krclass = {\begingroup\krfont}
\XeTeXinterchartoks \krclass 0 = {\endgroup}
\XeTeXinterchartoks \krclass 4095 = {\endgroup}
\XeTeXlinebreaklocale "ko"
\XeTeXlinebreakskip=0pt plus 1pt
\usepackage{amsmath,amssymb,amsthm,mathtools,bm}
\usepackage[italic,defaultmathsizes,noasterisk]{mathastext}
\emergencystretch=2.5em

\usepackage{graphicx,booktabs,tabularx,longtable,array,multirow,enumitem,xcolor}
\usepackage{caption}
\captionsetup{font=small,labelfont=bf}
\usepackage{tikz}
\usetikzlibrary{shapes.geometric,arrows.meta,positioning,fit,calc,backgrounds,decorations.pathreplacing}
\usepackage[most]{tcolorbox}
\usepackage{listings}
\usepackage{fancyhdr}
\usepackage[hidelinks]{hyperref}

\setlist{itemsep=1pt,topsep=3pt,parsep=0pt}
\setlength{\parskip}{3pt}
\renewcommand{\arraystretch}{1.12}

\lstdefinestyle{sql}{language=SQL,basicstyle=\ttfamily\small,keywordstyle=\bfseries,
  commentstyle=\itshape\color{black!60},columns=fullflexible,keepspaces=true,
  frame=l,framerule=0.6pt,rulecolor=\color{black!40},xleftmargin=8pt,framexleftmargin=4pt,
  aboveskip=5pt,belowskip=5pt,showstringspaces=false,breaklines=true,
  morekeywords={OVER,PARTITION,RANK,DENSE_RANK,ROW_NUMBER,NTILE,LAG,LEAD,ROWS,RANGE,PRECEDING,FOLLOWING,UNBOUNDED,CURRENT,ROW,NVL,NVL2,REFERENCES,CASCADE,MINUS,INTERSECT,EXISTS,MATERIALIZED,VIEW,REPLACE,RATIO_TO_REPORT,FIRST_VALUE,LAST_VALUE,INTERVAL,NULLS,LAST,FIRST,CHECK,CONSTRAINT,UNIQUE,DEFERRABLE,COMMIT,BEGIN,ROLLBACK,TRANSACTION,INDEX,DEFAULT,ANY,ALL,SOME,ROUND}}
\lstset{style=sql}

\tcbset{
 nb/.style={enhanced,breakable,colback=white,coltitle=black,fonttitle=\bfseries,
   left=3mm,right=3mm,top=1.6mm,bottom=1.6mm,before skip=9pt,after skip=9pt,parbox=false},
 framed/.style={nb,boxrule=0.5pt,colframe=black!70,arc=1.2mm,colbacktitle=white,
   titlerule=0.4pt,titlerule style={black!25},toptitle=1mm,bottomtitle=0.6mm},
 light/.style={framed,colframe=black!35},
 ruled/.style={nb,boxrule=0pt,leftrule=1.1pt,colframe=black!45,sharp corners,left=3.5mm,
   top=0.5mm,bottom=0.5mm,detach title,before upper={\tcbtitle\par\smallskip}},
}
\DeclareTColorBox{keybox}{o}{framed,leftrule=2.2pt,title={\IfValueTF{#1}{#1}{Key idea}}}
\DeclareTColorBox{formulabox}{o}{framed,IfValueTF={#1}{title={#1}}{}}
\DeclareTColorBox{workout}{o}{light,title={Example\IfValueT{#1}{\ \textmd{---}\ #1}}}
\DeclareTColorBox{qa}{m}{light,title={#1}}
\DeclareTColorBox{warnbox}{o}{ruled,leftrule=2.2pt,colframe=black,title={Caution\IfValueT{#1}{\ \textmd{---}\ #1}}}
\DeclareTColorBox{whybox}{o}{ruled,title={\IfValueTF{#1}{#1}{Why it works}}}
\DeclareTColorBox{prof}{o}{ruled,title={Lecture note\IfValueT{#1}{\ \textmd{---}\ #1}}}
\DeclareTColorBox{career}{o}{ruled,title={Career note\IfValueT{#1}{\ \textmd{---}\ #1}}}
\DeclareTColorBox{insight}{o}{ruled,title={In practice\IfValueT{#1}{\ \textmd{---}\ #1}}}

\newcommand{\Step}[1]{\par\smallskip\noindent\textbf{Step #1.}\ }
\newcommand{\Concept}[1]{\hfill{\small[개념: #1]}}
\newcommand{\Use}[1]{\noindent\textit{사용 개념:} #1\par}
\newcommand{\Ans}[1]{\par\smallskip\noindent\fbox{\parbox{\dimexpr\linewidth-2\fboxsep-2\fboxrule\relax}{\textbf{답}\ \ #1}}\par}
\newcommand{\Trap}[1]{\par\smallskip\noindent\textit{흔한 실수:} #1}
\newcommand{\Verify}[1]{\par\smallskip\noindent\textit{검산:} #1}
\newcommand{\Intuit}[1]{\par\smallskip\textbf{직관: #1}}
\newcommand{\pk}[1]{\underline{#1}}
\newcommand{\NULL}{\textsc{null}}
\newcommand{\sqlk}[1]{\texttt{#1}}
\newcommand{\Nul}{{\color{black!50}\textit{null}}}

% relational algebra symbols
\newcommand{\sel}{\sigma}
\newcommand{\proj}{\pi}
\newcommand{\join}{\bowtie}

\pagestyle{fancy}
\fancyhf{}
\fancyhead[R]{\small\itshape\leftmark}
\fancyfoot[C]{\small\thepage}
\renewcommand{\headrulewidth}{0pt}
\renewcommand{\chaptermark}[1]{\markboth{#1}{}}

\usepackage{titlesec}
\titleformat{\chapter}[display]{\bfseries\Large}{\chaptertitlename\ \thechapter}{6pt}{\LARGE}
\titlespacing*{\chapter}{0pt}{-10pt}{14pt}
\titleformat{\section}{\bfseries\large}{\thesection}{0.8em}{}
\titleformat{\subsection}{\bfseries\normalsize}{\thesubsection}{0.8em}{}

% ---------- cheat-sheet / problem-sheet helpers ----------
\newcounter{qn}
\newcommand{\Q}[2]{\par\smallskip\refstepcounter{qn}\noindent\textbf{Q\theqn.}\ {\footnotesize[#1]}\ #2\par}
\newcommand{\A}[1]{\vspace{1pt}\noindent\begin{tcolorbox}[ruled,leftrule=1.6pt,colframe=black!45,before skip=1pt,after skip=4pt,top=0.5mm,bottom=0.5mm]\small\textbf{답} #1\end{tcolorbox}}
\newcommand{\Lab}[3]{\item \textbf{#1}\\ \textit{해 본 것:} #2\\ \textit{\ensuremath{\Rightarrow} 개념:} #3}
\newcommand{\eg}{\textit{예)}\ }
\newcommand{\bi}{\textit{비유)}\ }
